\section{Random-phase encounter pairs and the exact $\tau^2$ coefficient}
\label{sec:random-phase-encounter}
The ensemble is $U=D_\theta M_N$ with $S=\begin{pmatrix}3&4\\2&3\end{pmatrix}$ and $D_\theta$ Haar on the $N$ position phases, $N$ odd. It is the maximal diagonal-phase ensemble. It is not the sixteen-phase ensemble of BGS-TRACE-126 and not the fixed kick. The canonical proof is
\path{mathematics/fixed-map/random_phase_encounter_pairs.md}.
Finite instances are certified by
\path{formalization/Bgs/LineCEncounter.lean}
in the namespace \texttt{Bgs.LineCEncounter}.
An independent referee engine is deposited at
\path{computations/verification/lineC/independent_replay/}.
BGS-TRACE-005 and BGS-BGS-001 remain open.

\begin{theorem}[Random-phase DFT identification, BGS-TRACE-143]
\label{thm:linec-dft}
Let $N$ be odd and $\Delta=\mathrm{diag}\,e_N(-3\cdot 8^{-1}x^2)$. Then $M_N=\Delta F_{1/4}\Delta$, and conjugation by $(Pf)(x)=f(2x)$ sends $F_{1/4}$ to the DFT. For every $\theta$, $D_\theta M_N$ is unitarily conjugate to $D_{\theta'}F$, and $\theta\mapsto\theta'$ preserves Haar measure. The eigenvalue law is that of Tanner's random-phase DFT star graph, which is the class-AI star graph of Gnutzmann--Seif (2004). For a unit cross coefficient the law depends only on its square class. At a prime $p\equiv 1\pmod 4$ a non-square class flips the sign of every half-integer term of $\mathbb E\lvert\operatorname{Tr} U^m\rvert^2$ and leaves every integer-order term unchanged. At $p=13$, $m=4$ the two values are $14033/2197\pm 8\sqrt{13}/2197$. The statement is for odd $N$ only.
\end{theorem}

\begin{theorem}[Gauss-sum formula, BGS-TRACE-144]
\label{thm:linec-gauss}
For odd $N$,
\[
\mathbb E\lvert\operatorname{Tr} U^m\rvert^2
=N^{-m}\sum_{\rho\in\Pi_m}W(\rho)\sum_{\sigma\in S_m}G_N(\rho,\sigma),
\]
with $G_N=\sum_z e_N(c\,Q_{\rho,\sigma}(z))$, $c=4^{-1}$, and $W(\rho)$ the product of the block weights whose exponential generating function is $\log\sum t^n/(n!)^2$. At a prime the sum is $p^{n-r/2}\varepsilon_p^r(\delta/p)$. The value is rational for every $p\equiv 3\pmod 4$, lies in $\mathbb Q(\sqrt p)$ for $p\equiv 1\pmod 4$, and has no $p^{-1/2}$ term. For $p>4^m$ the rank and the square class over $\mathbb F_p$ equal those over $\mathbb Q$. Lean decides the counting step at $N=7,11,13$ and $m=3$ only ($235/49$, $631/121$, $901/169$).
\end{theorem}

\begin{theorem}[Pair classification, BGS-TRACE-145]
\label{thm:linec-pairs}
At deficit $0$, $\mathrm{rank}_{\mathbb Q}(A_C-A_H)=2$ for a Hamiltonian cycle $H\neq C$ if and only if $H$ is a $2$-opt of $C$ ($m(m-3)/2$ cycles) or a swap of two vertices at cyclic distance at least $3$ ($m(m-5)/2$ cycles, $m\ge 5$). Each contributes a hyperbolic Gauss sum $N^{n-1}$ for every odd $N$. The diagonal is the dihedral group, $2m$ permutations for $m\ge 3$. At deficit $1$ and rank $0$ the pairs are the figure-eight words, with weight $w(2)=-1/2$. The swap class is absent from the three diagram families of Gnutzmann--Seif (2004). Lean decides the $m=6$ counts only. Ranks $1$ and $3$ are absent at deficit $0$ for $m\le 9$ in the note and for $m\le 10$ in the referee enumeration; that absence is not a general lemma.
\end{theorem}

\begin{theorem}[Exact $\tau^2$ coefficient, BGS-TRACE-146]
\label{thm:linec-tau2}
For $m\ge 3$ and every prime $p>4^m$,
$\mathbb E\lvert\operatorname{Tr}(DM_p)^m\rvert^2=\sum_{k=0}^{2m-1}a_{m,k}(p)\,p^{-k/2}$ exactly. One has $a_0=2m$ and $a_1=0$ for every $m$, and $a_2=-2m^2$ for every $m\ge 5$, with $-9$ at $m=3$ and $-24$ at $m=4$. This is the COE coefficient $-2\tau^2$. The $2$-opt pairs contribute $m^2(m-3)$, the swaps $m^2(m-5)$, and the figure eights $-2m^3+6m^2$. The remainder satisfies $\lvert\mathbb E-2m-a_2/p\rvert\le C_m p^{-3/2}$ with $C_m=m!\sum\lvert W\rvert$ ($3651$, $90921$, $3.1\cdot 10^6$, $1.4\cdot 10^8$ at $m=5,\ldots,8$). That explicit bound is informative only for $p>(C_m/2m^2)^2$. The statement is at fixed $m$. Gnutzmann--Seif (2004) derived $2\tau-2\tau^2$ diagrammatically for this ensemble in the limit $n,B\to\infty$. The coefficient $-2\tau^2$ is not new. What is new is the rigorous fixed-$m$, exact-in-$p$ identity and the complete pair classification, including swaps. The values at $m=1,2,3$ are BGS-TRACE-127 and are not re-claimed. Lean decides the integer identity for $5\le m\le 300$ and the hyperbolic counts at $p=5,7,11,13$.
\end{theorem}

\begin{theorem}[Exact expansions through $m=12$, BGS-TRACE-147, verification-pending]
\label{thm:linec-expansions}
For $p>4^m$ the expansion is known to all orders for $m\le 10$ and through $p^{-5/2}$ for $m=11,12$. There is no $p^{-3/2}$ term for $m\le 12$. The coefficient $c_2$ equals $2m^3$ at $m=7,10,11$ and differs at $m=5$ ($175+10(5/p)$), $6$ ($390$), $8$ ($1008$), $9$ ($1512$) and $12$ ($3480$ against $3456$). This is a fixed-$m$ deviation from the finite-$N$ COE value $2m^3$. It does not contradict Gnutzmann--Seif, whose result stops at $O(\tau^3)$. The $p^{-5/2}$ coefficients for the square class, at $p\equiv 1\pmod 4$, are $8,-72,128(2/p),200$ at $m=4,6,8,10$ and $0$ at $m=12$. An independent from-scratch engine reproduces every coefficient at $171$ primes. The claim stays verification-pending: the ranges are computational, and promotion is a separate status decision, as for BGS-RES-022. There is no Lean certificate of these expansions.
\end{theorem}

BGS-TRACE-148 remains open. It conjectures $c_2(m)=2m^3+O(m^2)$, that each trivial-character coefficient of $p^{-k}$ is the COE leading term times $(1+O(1/m))$, and that the half-integer and character terms vanish in the $\tau$-scaling limit. The guess $c_2=2m^3$ for all $m\ge 10$ is refuted at $m=12$. The $m\le 12$ table is consistent with the conjecture and does not prove it. Nothing here is a fixed-$\tau$ theorem, and nothing covers the fixed kick.
